\hypertarget{appendix-d-the-complete-python-grammar-ebnf}{%
\chapter{Appendix D: The Complete Python Grammar
(EBNF)}\label{appendix-d-the-complete-python-grammar-ebnf}}

This appendix provides the formal EBNF (Extended Backus-Naur Form)
grammar for Python 3.13. Understanding this grammar is the final step in
``Godhood,'' as it allows you to predict how any sequence of tokens will
be parsed by the PEG engine.

\hypertarget{d.1-notation}{%
\subsection{D.1 Notation}\label{d.1-notation}}

\begin{itemize}
\tightlist
\item
  \passthrough{\lstinline!?!} : Optional
\item
  \passthrough{\lstinline!*!} : 0 or more
\item
  \passthrough{\lstinline!+!} : 1 or more
\item
  \passthrough{\lstinline!|!} : Choice
\item
  \passthrough{\lstinline!()!} : Grouping
\end{itemize}

\hypertarget{d.2-the-core-grammar-abridged}{%
\subsection{D.2 The Core Grammar
(Abridged)}\label{d.2-the-core-grammar-abridged}}

\begin{lstlisting}
file: [statements] ENDMARKER
interactive: statement_newline

statements: statement+
statement: compound_stmt | simple_stmts

simple_stmts:
    | simple_stmt ';' [simple_stmts] NEWLINE
    | simple_stmt NEWLINE

simple_stmt:
    | assignment
    | type_alias
    | star_expressions
    | return_stmt
    | import_stmt
    | raise_stmt
    | 'pass'
    | del_stmt
    | yield_stmt
    | assert_stmt
    | 'break'
    | 'continue'
    | global_stmt
    | nonlocal_stmt

compound_stmt:
    | function_def
    | if_stmt
    | class_def
    | with_stmt
    | for_stmt
    | try_stmt
    | while_stmt
    | match_stmt

assignment:
    | NAME ':' expression ['=' annotated_rhs]
    | ('(' single_target ')' | single_subscript_attribute_target) ':' expression ['=' annotated_rhs]
    | (star_targets '=' )+ (yield_expr | star_expressions) [TYPE_COMMENT]
    | target _augassign_op (yield_expr | star_expressions)

if_stmt:
    | 'if' named_expression ':' block elif_stmt
    | 'if' named_expression ':' block [else_block]

elif_stmt:
    | 'elif' named_expression ':' block elif_stmt
    | 'elif' named_expression ':' block [else_block]

for_stmt:
    | [ASYNC] 'for' star_targets 'in' star_expressions ':' [TYPE_COMMENT] block [else_block]

while_stmt:
    | 'while' named_expression ':' block [else_block]

try_stmt:
    | 'try' ':' block finally_block
    | 'try' ':' block except_block+ [else_block] [finally_block]
    | 'try' ':' block except_star_block+ [else_block] [finally_block]

match_stmt:
    | "match" subject_expr ':' NEWLINE INDENT case_block+ DEDENT

case_block:
    | "case" patterns [guard] ':' block

# ... [Many pages of expressions, atoms, and literals] ...

expressions:
    | expression (',' expression )* [',']

expression:
    | conditional_expression
    | lambdef

conditional_expression:
    | disjunction 'if' disjunction 'else' expression
    | disjunction

disjunction:
    | conjunction ( 'or' conjunction )+
    | conjunction

conjunction:
    | inversion ( 'and' inversion )+
    | inversion

inversion:
    | 'not' inversion
    | comparison

comparison:
    | bitwise_or ( compare_op_bitwise_or_pair )+
    | bitwise_or

bitwise_or:
    | bitwise_or '|' bitwise_xor
    | bitwise_xor

bitwise_xor:
    | bitwise_xor '^' bitwise_and
    | bitwise_and

bitwise_and:
    | bitwise_and '&' shift_expr
    | shift_expr

shift_expr:
    | shift_expr '<<' sum
    | shift_expr '>>' sum
    | sum

sum:
    | sum '+' term
    | sum '-' term
    | term

term:
    | term '*' factor
    | term '/' factor
    | term '//' factor
    | term '%' factor
    | term '@' factor
    | factor

factor:
    | '+' factor
    | '-' factor
    | '~' factor
    | power

power:
    | await_primary '**' factor
    | await_primary

await_primary:
    | AWAIT primary
    | primary

primary:
    | primary '.' NAME
    | primary '(' [arguments] ')'
    | primary '[' slices ']'
    | atom

atom:
    | NAME
    | 'True'
    | 'False'
    | 'None'
    | strings
    | NUMBER
    | (tuple | list | dict | set)
    | '...'
\end{lstlisting}

\hypertarget{d.3-implications-of-peg}{%
\subsection{D.3 Implications of PEG}\label{d.3-implications-of-peg}}

Because Python uses a PEG (Parsing Expression Grammar) parser, the order
of rules in a choice (\passthrough{\lstinline!|!}) matters. The parser
tries the first option, and if it succeeds, it never looks at the
others. This eliminates ambiguity but requires careful ordering (e.g.,
matching longer keywords before shorter ones).

\begin{center}\rule{0.5\linewidth}{0.5pt}\end{center}

\textbf{THE END.}

\begin{center}\rule{0.5\linewidth}{0.5pt}\end{center}
